% =============================================================================
%  Un Vision Transformer desde cero en C con aceleración por GPU
%  Ezequiel González Tortajada
%
%  Compilar con:  tectonic memoria.tex
% =============================================================================
\documentclass[11pt,a4paper]{article}

\usepackage[spanish,es-noquoting,es-noshorthands,es-tabla]{babel}
\usepackage{lmodern}
\usepackage[T1]{fontenc}
\usepackage[margin=2.6cm]{geometry}
\usepackage{booktabs}
\usepackage{amsmath,amssymb}
\usepackage{microtype}
\usepackage{caption}
\usepackage{float}
\usepackage{xcolor}
\usepackage[htt]{hyphenat}
\usepackage{pgfplots}
\pgfplotsset{compat=1.18}
\emergencystretch=1.5em
\usepackage[hidelinks]{hyperref}

\captionsetup{font=small,labelfont=bf}
\setlength{\parskip}{0.4em}


% =============================================================================
\title{Un Vision Transformer desde cero en C con aceleración por GPU\\[0.3em]
       \large \textquestiondown{}Compensa frente a una arquitectura definida en Python?}
\author{Ezequiel González Tortajada}
\date{26 de septiembre de 2026}

\begin{document}
\maketitle

\begin{abstract}
\noindent
Este trabajo implementa desde cero, en C99 y sin librerías de aprendizaje
profundo, un \emph{Vision Transformer} completo: tensores, retropropagación
de cada operación, optimizador AdamW, aumentación de datos,
evaluación con \emph{ensemble} y \emph{test-time augmentation}, y dos backends de
GPU, uno para Metal (Apple Silicon) y otro para CUDA (NVIDIA), que ejecutan el
paso de entrenamiento completo en el dispositivo. Sobre el conjunto de datos
OrganSMNIST, con cortes de tomografía abdominal de $28\times28$ píxeles
clasificados en 11 órganos, un modelo de 0.80 millones de parámetros entrenado
sin pre-entrenamiento alcanza un 75.9\,\%
de precisión en test, y un \emph{ensemble} de dos modelos con TTA un 78.6\,\%.
El modelo individual queda 2.8~puntos por debajo de una ResNet-50 pre-entrenada
y ajustada a la misma resolución en PyTorch (78.7\,\%, 23.5 millones de
parámetros), a la que el \emph{ensemble} iguala; ambos quedan por debajo del
mejor modelo de referencia, un ViT-small pre-entrenado a $64\times64$
(82.2\,\%). La comparación muestra que el
rendimiento computacional se puede alcanzar con un esfuerzo acotado, pero que el
coste real de construir desde cero está en la verificación: el error más grave
del proyecto, un gradiente nulo en la atención, redujo la precisión al
51.6\,\% sin producir ningún fallo visible. Se concluye que, para obtener la
mejor precisión en un \emph{benchmark}, una arquitectura definida en Python es
claramente preferible; la implementación propia compensa cuando el objetivo es
comprender el modelo, desplegarlo sin dependencias o controlar los kernels.
\end{abstract}

\tableofcontents
\newpage

% =============================================================================
\section{Introducción}

Entrenar un \emph{Vision Transformer}~\cite{vit} con PyTorch~\cite{pytorch} y la
librería \texttt{timm}~\cite{timm} requiere unas decenas de líneas: la
arquitectura, la diferenciación automática, el optimizador y la ejecución en GPU
vienen resueltos, y se dispone de pesos pre-entrenados sobre millones de
imágenes. Este trabajo recorre el camino opuesto: construir el mismo tipo de
modelo desde cero en C, sin dependencias, incluyendo su ejecución en GPU.

La pregunta que guía la memoria es si ese esfuerzo compensa. Para responderla
con datos, se entrena en el mismo portátil y con el mismo protocolo una
referencia en Python: 13 modelos en PyTorch y scikit-learn, desde modelos
clásicos hasta redes pre-entrenadas con \emph{fine-tuning}, entre ellas un
ViT-small.

Las contribuciones son:

\begin{itemize}
  \item Un ViT entrenable de extremo a extremo en C99, con retropropagación
        verificada por diferencias finitas.
  \item Un diseño de backend de dispositivo con memoria compartida y una sola
        sincronización por paso de entrenamiento, implementado para Metal y
        para CUDA.
  \item Kernels de atención en tres variantes, incluida una por bloques de tipo
        FlashAttention~\cite{flashattention} para secuencias largas.
  \item Una comparación cuantitativa de precisión, coste computacional y
        esfuerzo de desarrollo frente a modelos definidos en Python.
\end{itemize}

% =============================================================================
\section{Datos y protocolo}

\subsection{OrganSMNIST}

OrganSMNIST pertenece a la colección MedMNIST v2~\cite{medmnistv2}. Contiene
cortes sagitales de tomografía computarizada abdominal de $28\times28$ píxeles
en escala de grises, etiquetados en 11 órganos, con una partición oficial de
13\,932 imágenes de entrenamiento, 2\,452 de validación y 8\,827 de test. La
tabla~\ref{tab:oficial} recoge los resultados publicados por los autores de la
colección.

\begin{table}[H]
  \centering\small
  \begin{tabular}{lrr}
    \toprule
    método & AUC & precisión \\
    \midrule
    ResNet-18 (28)       & 0.972 & 78.2\,\% \\
    ResNet-18 (224)      & 0.974 & 77.8\,\% \\
    ResNet-50 (28)       & 0.972 & 77.0\,\% \\
    ResNet-50 (224)      & 0.975 & 78.5\,\% \\
    auto-sklearn         & 0.945 & 67.2\,\% \\
    AutoKeras            & 0.974 & 81.3\,\% \\
    Google AutoML Vision & 0.964 & 74.9\,\% \\
    \bottomrule
  \end{tabular}
  \caption{Resultados oficiales de MedMNIST v2 sobre el test de OrganSMNIST.}
  \label{tab:oficial}
\end{table}

\subsection{Protocolo}

Todos los modelos, propios y de referencia, siguen el mismo protocolo, lo que
hace comparables las cifras:

\begin{itemize}
  \item \emph{Particiones.} Entrenamiento ajusta los pesos; validación
        selecciona el checkpoint, la arquitectura y la configuración de
        evaluación; test se consulta solo para informar.
  \item \emph{Preproceso.} Los píxeles se normalizan con la media y la
        desviación típica del conjunto de entrenamiento (0.495 y 0.283).
  \item \emph{Métricas.} Precisión, precisión equilibrada (media del
        \emph{recall} por clase), macro-F1 y matriz de confusión.
\end{itemize}

\subsection{Dos versiones de las imágenes}
\label{sec:versiones}

OrganSMNIST circula en dos formatos: el fichero oficial, con las imágenes ya
reducidas a $28\times28$ por los autores de la colección, y copias en
repositorios de terceros con las imágenes en alta resolución. Parte del
desarrollo se hizo con estas últimas, reducidas a $28\times28$ con la
interpolación de la librería Pillow. Las particiones y las etiquetas coinciden,
pero los píxeles no: la desviación típica del conjunto de entrenamiento es 0.262
en esa versión y 0.283 en la oficial.

La diferencia no es inocua. Con la misma configuración, la versión reducida con
Pillow alcanza un 78.4 y un 78.6\,\% en test en dos ejecuciones, y la oficial un
75.9 y un 76.2\,\% con dos semillas. Todas las cifras principales de esta
memoria usan la versión oficial, que es la de los modelos de referencia y la de
la tabla~\ref{tab:oficial}. Solo el estudio de arquitecturas de la
sección~\ref{sec:arquitecturas}, que compara variantes entre sí, usa la otra
versión.

% =============================================================================
\section{Implementación}

\subsection{Modelo}

El modelo es un ViT con normalización previa (\emph{Pre-LN})~\cite{preln}. Cada
imagen se divide en parches no solapados de $P\times P$ píxeles, que se
proyectan a dimensión $d$; se antepone un token \texttt{[CLS]} y se suman
posiciones aprendibles. Cada bloque aplica
\begin{align*}
  h &= x + \mathrm{MHA}(\mathrm{LN}(x)), \\
  y &= h + W_2\,\mathrm{GELU}(W_1\,\mathrm{LN}(h)),
\end{align*}
con atención multi-cabeza bidireccional y un MLP de expansión 4. Tras el último
bloque, una LayerNorm y una capa lineal sobre \texttt{[CLS]} producen los
\emph{logits}. La configuración principal ($P=4$, $d=128$, 8 cabezas, 4 bloques)
tiene 50 tokens por imagen y 803\,467 parámetros.

\subsection{Entrenamiento}

No hay diferenciación automática: cada operación tiene un \emph{forward} que
guarda en una caché las activaciones necesarias y un \emph{backward} escrito a
mano que acumula gradientes. El entrenamiento usa AdamW~\cite{adamw} con
\emph{weight decay} desacoplado, \emph{clipping} global por norma, \emph{label
smoothing}~\cite{labelsmoothing}, \emph{linear warmup} y \emph{cosine decay}. La
aumentación de datos se limita a transformaciones que conservan la orientación
anatómica (traslación, rotación, escala, brillo, contraste, ruido y
\emph{random erasing}). Se excluyen los volteos, que invierten la orientación
de la imagen, siempre la misma en los datos.

La evaluación admite \emph{ensembles} de checkpoints, que promedian
probabilidades y pueden mezclar arquitecturas, y \emph{test-time augmentation}
con 5 o 9 vistas desplazadas un píxel.

\subsection{Verificación}

Como no hay diferenciación automática, cada gradiente es una fuente potencial
de error silencioso. La verificación se apoya en dos mecanismos:

\begin{itemize}
  \item \emph{Diferencias finitas.} Se comparan los gradientes analíticos con
        $\frac{f(\theta+\epsilon)-f(\theta-\epsilon)}{2\epsilon}$ en las
        operaciones y en el bloque completo.
  \item \emph{CPU como referencia.} La implementación en C es la referencia
        numérica. Un test entrena el mismo modelo en CPU y en GPU y exige que
        coincidan la pérdida, los gradientes de todos los parámetros y la
        trayectoria de entrenamiento, en tres configuraciones que cubren las
        tres variantes del kernel de atención.
\end{itemize}

\subsection{Backends de GPU}

Ambos backends siguen el mismo diseño, resumido en la
tabla~\ref{tab:backends}:

\begin{enumerate}
  \item \emph{Memoria compartida.} Al activar el backend, cada tensor se reserva
        en memoria accesible desde CPU y GPU. La CPU solo accede a ella tras
        sincronizar.
  \item \emph{Per-operation dispatch.} Cada operación de la librería consulta
        el backend instalado; si todos sus tensores residen en el dispositivo,
        encola el kernel sin esperar, y si no, sincroniza y ejecuta la versión
        en C.
  \item \emph{Una sincronización por paso.} El paso de entrenamiento completo,
        desde la extracción de parches hasta AdamW, se encola entero y solo se
        sincroniza al leer la pérdida.
\end{enumerate}

\begin{table}[H]
  \centering\small
  \begin{tabular}{lll}
    \toprule
    & Metal & CUDA \\
    \midrule
    memoria compartida    & \texttt{MTLBuffer} compartido & \texttt{cudaMallocManaged} \\
    multiplicación de matrices & \texttt{MPSMatrixMultiplication} & cuBLAS \texttt{cublasSgemm} \\
    cola de trabajo       & \emph{command buffer} pendiente & \texttt{cudaStream\_t} \\
    reducciones           & SIMD-groups de 32 hilos & \emph{warps} \\
    parámetros agrupados  & tabla de \texttt{gpuAddress} & tabla de punteros \\
    \bottomrule
  \end{tabular}
  \caption{Correspondencia entre los dos backends.}
  \label{tab:backends}
\end{table}

La atención elige entre tres kernels. Si la secuencia completa cabe en memoria
compartida ($P=4$, 50 tokens), un grupo de 256 hilos por par (muestra, cabeza)
calcula todas las fases en paralelo; el \emph{backward} usa
$D_i = dO_i\cdot O_i$ para no almacenar $dP$~\cite{flashattention}. Si no cabe
($P=2$, 197 tokens), un kernel por bloques de 32 consultas recorre las claves
por bloques con \emph{softmax} en línea. Para cabezas de dimensión mayor que 64
queda un kernel genérico.

El backend CUDA replica el de Metal kernel a kernel y se ha validado en una GPU
NVIDIA L40S. El test de comparación con la CPU pasa en las tres configuraciones
de la atención. Además, con los mismos datos y la misma semilla, las dos
primeras épocas de entrenamiento (870 pasos) producen en la L40S y en el M3 la
misma pérdida y la misma precisión de validación hasta la sexta cifra decimal:
dos GPU de fabricantes distintos, con kernels escritos por separado, reproducen
el mismo entrenamiento.

% =============================================================================
\section{Rendimiento computacional}

Todas las medidas se toman en un MacBook con chip M3 (8 núcleos de CPU y GPU
integrada), con \emph{batch} de 32 y la configuración principal. La
tabla~\ref{tab:rendimiento} resume la evolución del tiempo por paso de
entrenamiento.

\begin{table}[H]
  \centering\small
  \begin{tabular}{lrr}
    \toprule
    versión & ms por paso & muestras por segundo \\
    \midrule
    CPU, multiplicación de matrices ingenua       & 3\,967 & 8 \\
    CPU en serie + multiplicaciones en GPU (MPS)  & 150    & 213 \\
    + paralelismo en CPU (8 núcleos)              & 58.5   & 547 \\
    + multiplicaciones sin transponer en CPU      & 50.0   & 640 \\
    + operaciones de los bloques en GPU           & 28.6   & 1\,121 \\
    + paso completo en GPU, una sincronización    & 24.6   & 1\,300 \\
    + kernels optimizados                         & 15.5   & 2\,065 \\
    + \emph{clipping} y AdamW agrupados           & 15.0   & 2\,133 \\
    \bottomrule
  \end{tabular}
  \caption{Tiempo por paso de entrenamiento ($P=4$, $d=128$, 4 bloques).}
  \label{tab:rendimiento}
\end{table}

\begin{figure}[H]
  \centering
  \begin{tikzpicture}
    \begin{axis}[
      width=0.85\linewidth, height=5.2cm,
      ybar, bar width=14pt,
      ylabel={ms por paso},
      symbolic x coords={serie,paralelo,GEMM,bloques,paso,kernels,agrupado},
      xtick=data, xticklabel style={font=\small},
      nodes near coords, nodes near coords style={font=\scriptsize},
      ymin=0, ymax=170, enlarge x limits=0.08,
    ]
      \addplot coordinates {(serie,150) (paralelo,58.5) (GEMM,50) (bloques,28.6)
                            (paso,24.6) (kernels,15.5) (agrupado,15.0)};
    \end{axis}
  \end{tikzpicture}
  \caption{Evolución del tiempo por paso a partir de la primera versión con GPU.}
  \label{fig:rendimiento}
\end{figure}

Tres observaciones se desprenden del perfilado:

\begin{itemize}
  \item \emph{La GPU solo acelera si el trabajo reside en ella.} Enviar a la GPU
        únicamente las multiplicaciones de matrices, copiando operandos y
        esperando cada una, dejaba la GPU al 18\,\% de uso: el tiempo se iba en
        la CPU y en las sincronizaciones. El salto grande llega al mover todas
        las operaciones y encolar el paso completo.
  \item \emph{Los kernels ingenuos desperdician la GPU.} Las
        reducciones por columnas (gradientes de sesgos, $\gamma$ y $\beta$) con
        un hilo por columna recorrían 1\,600 filas en serie. Reescritas con
        bloques de $32\times32$ hilos pasaron de 4.1 a 0.3~ms por paso; la
        LayerNorm, con un SIMD-group por fila, de 3.3 a 0.3~ms.
  \item \emph{Un primer intento de optimización empeoró el resultado.} Cargar la
        atención en memoria compartida con un hilo por consulta dejaba 50 de
        256 hilos activos y multiplicó por 2.7 el tiempo del \emph{backward};
        repartir cada fase entre todos los hilos lo redujo a 0.65~ms por
        llamada.
\end{itemize}

Con parches de $2\times2$, la secuencia pasa de 50 a 197 tokens. El kernel por
bloques reduce el paso de 152 a 101~ms, pero cada época sigue costando unos
70~s frente a los 8~s de la configuración principal.

% =============================================================================
\section{Resultados}

\subsection{Resultado principal}

La tabla~\ref{tab:principal} recoge los resultados de la configuración
principal con la versión oficial de los datos. Se entrenan dos modelos con
semillas distintas y se combinan en un \emph{ensemble}; el número de vistas de
TTA se elige por la precisión de validación.

\begin{table}[H]
  \centering\small
  \begin{tabular}{lrrrr}
    \toprule
    modelo & validación & test & equilibrada & macro-F1 \\
    \midrule
    semilla 42                         & 88.1\,\% & 75.9\,\% & 0.718 & 0.721 \\
    semilla 7                          & 87.2\,\% & 76.2\,\% & 0.724 & 0.719 \\
    \emph{ensemble}, sin TTA            & 89.0\,\% & 77.4\,\% & 0.735 & 0.742 \\
    \emph{ensemble}, TTA de 5 vistas    & 89.8\,\% & 78.4\,\% & 0.745 & 0.752 \\
    \emph{ensemble}, TTA de 9 vistas    & \textbf{90.3\,\%} & \textbf{78.6\,\%} & \textbf{0.745} & \textbf{0.754} \\
    \bottomrule
  \end{tabular}
  \caption{Resultados de la configuración principal ($P=4$, $d=128$, 4 bloques)
           sobre la versión oficial de los datos.}
  \label{tab:principal}
\end{table}

Las dos semillas difieren en 0.3~puntos en test. Promediar ambos modelos aporta
1.4~puntos, y la TTA otros 1.2.

\subsection{Un error silencioso}

Durante el desarrollo, con la versión de los datos reducida con Pillow, el
\emph{backward} de la atención
enrutaba el gradiente de la salida en sentido contrario, de modo que las
proyecciones de consultas, claves y valores nunca recibían gradiente y
conservaban su inicialización aleatoria. El programa entrenaba, la pérdida
bajaba y la precisión alcanzaba un 51.6\,\%: un resultado plausible para un
modelo pequeño. Solo una comprobación por diferencias finitas sobre los pesos
de la atención, que el test existente no incluía, reveló que el gradiente era
exactamente cero. Tras corregirlo, junto con la normalización de la entrada y
una inicialización aleatoria real, la precisión en validación llegó al 81\,\%
en cinco épocas.

\subsection{Tamaño de la arquitectura}
\label{sec:arquitecturas}

La tabla~\ref{tab:arquitecturas} compara variantes entrenadas con los mismos
hiperparámetros sobre la versión de los datos reducida con Pillow
(sección~\ref{sec:versiones}). Sirve para comparar las variantes entre sí, no
con el resto de cifras de la memoria. Se seleccionan por la pérdida de
validación.

\begin{table}[H]
  \centering\small
  \begin{tabular}{lrrrr}
    \toprule
    configuración & parámetros & pérdida val. & precisión val. & test \\
    \midrule
    $P=4$, $d=128$, 4 bloques & 0.80\,M & \textbf{0.311} & \textbf{91.0\,\%} & 78.4\,\% \\
    $P=4$, $d=192$, 6 bloques & 2.68\,M & 0.363 & 89.0\,\% & 77.4\,\% \\
    $P=2$, $d=128$, 4 bloques & 0.82\,M & 0.347 & 89.1\,\% & 77.8\,\% \\
    \bottomrule
  \end{tabular}
  \caption{Variantes de arquitectura, con la versión de los datos reducida con
           Pillow.}
  \label{tab:arquitecturas}
\end{table}

Ninguna variante mejora la configuración de partida. Con 13\,932 imágenes y sin
pre-entrenamiento, aumentar la profundidad, la anchura o la resolución de los
parches no aporta: el cuello de botella no está en la capacidad del modelo.

\subsection{El límite de la lateralidad}

Casi un tercio de los errores del mejor modelo individual (el 32\,\%) se
concentra en dos pares: fémur izquierdo y derecho, y riñón izquierdo y derecho. Para medir cuánta
información contienen las imágenes sobre la lateralidad se clasifica cada par
con el criterio más simple posible, la distancia a la imagen media de cada
clase:

\begin{table}[H]
  \centering\small
  \begin{tabular}{lrrr}
    \toprule
    par & media más cercana (val.) & media más cercana (test) & ViT, dentro del par (test) \\
    \midrule
    fémur & 62.3\,\% & 51.5\,\% & 50.4\,\% \\
    riñón & 75.6\,\% & 61.8\,\% & 72.4\,\% \\
    \bottomrule
  \end{tabular}
  \caption{Acierto al separar el órgano izquierdo del derecho en cada par.}
  \label{tab:lateralidad}
\end{table}

En test, los fémures se separan al azar con cualquier criterio, y el ViT no lo
mejora. En los riñones queda algo de señal, que el ViT aprovecha. La
validación vuelve a ser más optimista que el test, como ocurre con todos los
modelos de esta memoria. La explicación es anatómica: en un corte sagital, la
lateralidad depende de la profundidad a la que se tomó el corte, no del
contenido de la imagen.

% =============================================================================
\section{Comparación con arquitecturas definidas en Python}

\subsection{Precisión y coste}

La tabla~\ref{tab:comparacion} reúne los principales modelos de referencia en
PyTorch y los de este trabajo. Todos se entrenaron en el mismo portátil,
PyTorch con su backend MPS y el ViT en C con el backend Metal propio.

\begin{table}[H]
  \centering\small
  \setlength{\tabcolsep}{4pt}
  \begin{tabular}{llcrrrrr}
    \toprule
    modelo & origen & pre-entr. & resol. & parámetros & test & equil. & tiempo \\
    \midrule
    perceptrón multicapa & PyTorch & no & 28 & 0.67\,M & 59.1\,\% & 0.558 & 28\,s \\
    LeNet                & PyTorch & no & 28 & 0.06\,M & 69.4\,\% & 0.644 & 40\,s \\
    ResNet-50, completo  & PyTorch & sí & 28 & 23.5\,M & 78.7\,\% & 0.746 & 407\,s \\
    ResNet-50            & PyTorch & sí & 64 & 23.5\,M & 80.9\,\% & 0.763 & 744\,s \\
    ViT-small            & PyTorch & sí & 64 & 21.4\,M & \textbf{82.2\,\%} & \textbf{0.784} & 1\,065\,s \\
    \midrule
    ViT desde cero       & C       & no & 28 & 0.80\,M & 75.9\,\% & 0.718 & 566\,s \\
    \emph{ensemble} $\times 2$ + TTA & C   & no & 28 & 1.61\,M & 78.6\,\% & 0.745 & 1\,208\,s \\
    \bottomrule
  \end{tabular}
  \caption{Modelos de referencia en PyTorch y modelos de este trabajo sobre el
           test de OrganSMNIST. El tiempo es el de entrenamiento: 12 épocas para los
           modelos pre-entrenados, 30 para las redes desde cero de PyTorch y 50
           para el ViT en C.}
  \label{tab:comparacion}
\end{table}

Cuatro lecturas:

\begin{itemize}
  \item \emph{El modelo individual queda algo por debajo; el ensemble, a la
        par.} Entrenado desde cero a $28\times28$, el ViT en C queda 2.8~puntos
        por debajo de la ResNet-50 pre-entrenada a la misma resolución y
        2.3~puntos por debajo de la ResNet-18 oficial (tabla~\ref{tab:oficial}),
        con 29 y 14 veces menos parámetros. El \emph{ensemble} con TTA alcanza
        el 78.6\,\%, al nivel de ambas. Frente a la red de PyTorch entrenada
        desde cero, LeNet, el modelo individual mejora 6.5~puntos.
  \item \emph{La diferencia con el mejor modelo la marca el ecosistema, no el
        lenguaje.} Los 6.3~puntos que separan el ViT en C del ViT-small de
        PyTorch (3.6 con el \emph{ensemble}) se deben sobre todo a dos cosas que
        el lenguaje no aporta por sí mismo: el pre-entrenamiento sobre
        ImageNet-21k y la resolución de $64\times64$.
        La sección anterior muestra que ampliar el modelo sin pre-entrenar no
        cierra esa distancia.
  \item \emph{El coste computacional es comparable.} Tras la optimización, una
        época cuesta entre 8 y 13~s según la ejecución, y las 50 épocas, unos
        10~minutos. Los tiempos no son directamente comparables, porque cambian
        el número de épocas, la resolución y el tamaño de los modelos, pero en
        el mismo portátil el ViT en C tarda algo más que el \emph{fine-tuning}
        de la ResNet-50 a $28\times28$ y la mitad que el del ViT-small.
  \item \emph{Las diferencias pequeñas no son significativas.} En la referencia,
        dos ejecuciones de la misma ResNet-50 difirieron en 1.1~puntos; aquí,
        dos semillas del ViT difieren en 0.3. Diferencias por debajo de un punto
        no permiten ordenar modelos.
\end{itemize}

\subsection{Esfuerzo de desarrollo}

La tabla~\ref{tab:esfuerzo} compara el tamaño del código, sin contar líneas en
blanco.

\begin{table}[H]
  \centering\small
  \begin{tabular}{lr}
    \toprule
    componente & líneas \\
    \midrule
    núcleo en C (tensores, operaciones, ViT, aumentación, datos) & 3\,940 \\
    backend Metal (Objective-C++ y kernels Metal) & 2\,673 \\
    backend CUDA & 1\,980 \\
    programa de entrenamiento y evaluación & 1\,021 \\
    tests & 1\,216 \\
    \midrule
    total de la implementación propia (un ViT y dos backends de GPU) & 10\,830 \\
    referencia en PyTorch y scikit-learn (13 modelos) & 899 \\
    \bottomrule
  \end{tabular}
  \caption{Tamaño del código de cada enfoque.}
  \label{tab:esfuerzo}
\end{table}

La referencia entrena 13 modelos, de la regresión logística al ViT-small, con
un orden de magnitud menos de código que la implementación propia de un único
modelo. Pero la diferencia importante no es de volumen, sino de tipo de error. Los fallos encontrados durante el desarrollo fueron:

\begin{itemize}
  \item el gradiente nulo en la atención, que no producía ningún fallo visible y
        dejaba el modelo en un 51.6\,\% de precisión;
  \item una inicialización pseudoaleatoria con estructura de diente de
        sierra, que no era aleatoria;
  \item un planificador coseno que llegaba a una tasa de aprendizaje exacta de
        cero en la última época y abortaba el entrenamiento antes de evaluar;
  \item un objeto liberado por la gestión automática de memoria de
        Objective-C al reordenar una caché, que cerraba el programa;
  \item una variable ocultada en el backend CUDA, detectada por análisis
        estático antes de ejecutarlo;
  \item dos versiones de las imágenes con píxeles distintos, que desplazaban
        los resultados 2.5~puntos sin que ningún test pudiera detectarlo.
\end{itemize}

Con PyTorch, la diferenciación automática elimina el primero, la inicialización
y los planificadores de la librería eliminan el segundo y el tercero, y la
gestión de memoria del \emph{framework} elimina el cuarto. Los dos últimos no
dependen del lenguaje: pueden ocurrir con cualquier herramienta. Los fallos de un
\emph{framework} maduro suelen ser visibles; los de una implementación propia
pueden producir un modelo que entrena, converge y da resultados plausibles
siendo incorrecto. Por eso una implementación propia exige infraestructura de
verificación: diferencias finitas, una referencia numérica independiente y tests
que comparen dispositivos.

% =============================================================================
\section{Conclusiones}

\subsection{\textquestiondown{}Compensa?}

Depende del objetivo.

\emph{Para obtener la mejor precisión en un problema concreto, no compensa.}
Una arquitectura definida en Python ofrece, con un orden de magnitud menos de
código, modelos pre-entrenados que aportan hasta 6~puntos sobre este problema,
diferenciación automática que elimina la clase de error más grave encontrada, y
un ciclo de experimentación mucho más corto: la referencia cubre 13 modelos con
menos de mil líneas.

\emph{Compensa cuando el objetivo es otro:}

\begin{itemize}
  \item \emph{Comprender el modelo.} Escribir cada retropropagación obliga a
        entender dónde fluye cada gradiente; el error de la atención solo se
        explica entendiendo la operación.
  \item \emph{Desplegar sin dependencias.} El modelo final ocupa 3.2~MB de
        pesos, 0.8~millones de parámetros, y se ejecuta con una librería en
        C99 sin intérprete de Python ni \emph{runtime} de aprendizaje profundo.
  \item \emph{Controlar el cómputo.} Decidir cuándo se sincroniza la GPU, cómo
        se reparten los hilos de un kernel o cómo se agrupan los parámetros
        permitió reducir el paso de entrenamiento de 150 a 15~ms, y ese
        conocimiento se traslada a cualquier \emph{framework}.
\end{itemize}

El resultado más relevante es que la parte que parecía más difícil, el
rendimiento en GPU, resultó abordable y medible paso a paso. Lo costoso fue la
corrección: asegurarse de que cada gradiente es el correcto.

\subsection{Recomendación}

Un enfoque mixto aprovecha lo mejor de ambos: prototipar y seleccionar el
modelo en Python, con pesos pre-entrenados; usar el modelo de Python como
referencia numérica; y portar a C solo la inferencia, o el entrenamiento cuando
el despliegue lo requiera, verificando capa a capa contra esa referencia.

\subsection{Trabajo futuro}

\begin{itemize}
  \item Medir el rendimiento del backend CUDA en GPU dedicadas y compararlo con
        PyTorch en el mismo hardware.
  \item Cargar pesos pre-entrenados de un ViT de PyTorch en la implementación en
        C, para separar el efecto del pre-entrenamiento del de la
        implementación.
  \item Entrenar a $64\times64$ con la atención por bloques.
  \item Precisión reducida (fp16 o bf16) en las multiplicaciones de matrices.
\end{itemize}

% =============================================================================
\begin{thebibliography}{9}

\bibitem{medmnistv2}
J.~Yang, R.~Shi, D.~Wei, Z.~Liu, L.~Zhao, B.~Ke, H.~Pfister y B.~Ni.
\newblock MedMNIST v2 - A large-scale lightweight benchmark for 2D and 3D
  biomedical image classification.
\newblock \emph{Scientific Data}, 10:41, 2023.
\newblock \url{https://doi.org/10.1038/s41597-022-01721-8}.

\bibitem{vit}
A.~Dosovitskiy \emph{et al.}
\newblock An image is worth 16x16 words: Transformers for image recognition at
  scale.
\newblock En \emph{ICLR}, 2021.
\newblock \url{https://arxiv.org/abs/2010.11929}.

\bibitem{preln}
R.~Xiong \emph{et al.}
\newblock On layer normalization in the Transformer architecture.
\newblock En \emph{ICML}, 2020.
\newblock \url{https://arxiv.org/abs/2002.04745}.

\bibitem{adamw}
I.~Loshchilov y F.~Hutter.
\newblock Decoupled weight decay regularization.
\newblock En \emph{ICLR}, 2019.
\newblock \url{https://arxiv.org/abs/1711.05101}.

\bibitem{labelsmoothing}
C.~Szegedy, V.~Vanhoucke, S.~Ioffe, J.~Shlens y Z.~Wojna.
\newblock Rethinking the Inception architecture for computer vision.
\newblock En \emph{CVPR}, 2016.
\newblock \url{https://doi.org/10.1109/CVPR.2016.308}.

\bibitem{flashattention}
T.~Dao, D.~Y.~Fu, S.~Ermon, A.~Rudra y C.~Ré.
\newblock FlashAttention: Fast and memory-efficient exact attention with
  IO-awareness.
\newblock En \emph{NeurIPS}, 2022.
\newblock \url{https://arxiv.org/abs/2205.14135}.

\bibitem{pytorch}
A.~Paszke \emph{et al.}
\newblock PyTorch: An imperative style, high-performance deep learning library.
\newblock En \emph{NeurIPS}, 2019.
\newblock \url{https://arxiv.org/abs/1912.01703}.

\bibitem{timm}
R.~Wightman.
\newblock PyTorch Image Models, 2019.
\newblock \url{https://github.com/huggingface/pytorch-image-models}.

\end{thebibliography}

\end{document}
